\section{Random Separating Hyperplane ($\rsh$) Theorem}
\label{sec:rsh}
In this section, we prove $\rsh$ for polytopes: If a point $p$  point is at a distance from a polytope $K$, then, the Gaussian measure of the set of {\em well-separating} hyperplanes has a positive lower bound depending on the number of vertices in $K$. More specifically, we show:

\begin{theorem}
\label{thm:sephp}
Suppose $K$ is a polytope in $\Re^d$ with $k$ vertices and diameter $\dia(K)$. Suppose  $a$ is a point in ${\bf R}^d$ and $\delta\in (0,1]$ with
\begin{align}
\label{eq:condhp}
\min_{y\in K} |a-y|\geq\delta \dia(K).
\end{align}
Let $V$ be an $m$-dimensional subspace containing $\lspan(K\cup \{ a\})$ and let $u$ be a random vector drawn from the normal distribution $N(0, I_m)$ in $V$. 
Then,
$$\prob_u\left[ (u\cdot a-\max_{y\in K}u\cdot y) \geq  |u| \cdot \delta \dia(K) \cdot \frac{\sqrt{\log k}}{\sqrt{\log k} + 4 \delta \sqrt{m}} \right]\geq
            \frac{1}{40}    k^{-10/\delta^2}.$$
\end{theorem}          

% \noindent
% {\bf Remark:}  Such a theorem does not hold for general convex sets.
% For example,  consider the Euclidean space $\Re^d$ and the $d-2$ dimensional sphere $K:=\{x \in \Re^d: |x|\leq 1, x_1=0\}$. Observe that the diameter $\dia(K)$ of $K$ is 1. 
% Let $a$ denote the point $(1,0,0,\ldots ,0)$. It is at distance 1 from $K$, and so the condition~\eqref{eq:condhp} is satisfied if we set $\rho$ to $1/2$.  But, with high probability,  a random unit length vector $u$ has $u\cdot a\approx \frac{c}{\sqrt d}$ for some constant $c$,  whereas the maximum of  $v\cdot x$ over all points $x$ in the sphere $X$ is about 1. Therefore the probability of the event in the statement of~\Cref{thm:sephp} is exponentially small in $d$. 

\begin{proof}
 Let $\zeta_1,\zeta_2,\ldots ,\zeta_k$ be the vertices of $K$. Let $b$ be the closest point in $K$ to $a$, and define 
$$w=\frac{a-b}{|a-b|}.$$
Then by standard Convex Geoemtry arguments, we have for all points $y \in K$:
\begin{equation}\label{101}
w\cdot y\leq w\cdot b.
\end{equation}
Let $u$ be as in the statement of the theorem. 
We can write $u$ as 
$$u=\lambda w+z, \mbox{ where }z\perp w.$$

In order to prove the theorem, we first express $(u \cdot a  - \max_{y \in K} u \cdot y)$ as 
\begin{align}
\label{eq:step1}
u \cdot (a-b) - \max_{y \in K} u \cdot (y-b)  =  u \cdot (a-b) - \max_{\ell = 1, \ldots, k} u \cdot (\zeta_\ell-b)
\end{align}

It suffices to show that with reasonably high probability, there is a lower bound on $u \cdot (a-b)$ and an upper bound on $u \cdot (\zeta_\ell-b)$ for all $\ell$. Observe that 
\begin{align}
\label{eq:step2}
u \cdot (\zeta_\ell-b) = \lambda w \cdot (\zeta_\ell-b) + z \cdot (\zeta_\ell-b) \stackrel{\eqref{101}}{\leq} z \cdot (\zeta_\ell-b).
\end{align}
Therefore, it suffices to upper bound $z \cdot (\zeta_\ell - b)$ for all vertices $\zeta_\ell$. We now define a sequence of events and bound the probability of each of these events. 

\begin{fact}
\label{fact:1}
Let $\cE_0$ be the event $[|z| \leq 4 \sqrt{m}]$. Then $\Pr[\neg \cE_0] \leq 1/20. $
\end{fact}
\begin{proof}
The coordinates of $z$ are independent $N(0,1)$ random variables, and since $z$ lies in an $(m-1)$-dimensional subspace, the result follows from standard concentration bounds for sum of squares of Gaussian random variables. 
\end{proof}

\begin{fact}
\label{fact:2}
For $\ell=1,2,\ldots , k$, define the  event
${\cal E}_\ell$ as $[\left| z\cdot (\zeta_\ell-b)\right|\leq 2\sqrt{\ln k} \dia(K)].$ Then $\Pr[\neg \cE_\ell] \leq \frac{1}{4k^2}.$
\end{fact}
\begin{proof}
Observe that $z \cdot (\zeta_\ell - b)$ is a 1-dimensional Gaussian random variable with zero mean and  with variance $|\zeta_\ell - b|^2 \leq \dia(K)^2$. The result now follows the tail bounds for a 1-dimensional Gaussian random variable. 
\end{proof}

\begin{fact}
\label{fact:3}
Let
${\cal E}_{k+1}$ denote the event that $\lambda\geq 3\sqrt{\ln k}/\delta.$ Then $\Pr[\neg \cE_{k+1}] \geq  \frac{1}{10} k^{-10/\delta^2}.$
\end{fact}
\begin{proof}
Since $\lambda$ is a 1-dimensional Gaussian random variable, the result follows from tail bounds for normal distribution. 
\end{proof}

\begin{corollary}
\label{cor:event}
$\prob\left[ \wedge_{\ell=0}^{k+1} {\cal E}_\ell\right] \geq \frac{1}{40}    k^{-10/\delta^2}.$
\end{corollary}

\begin{proof}
Applying union bound to~\Cref{fact:1} and~\Cref{fact:2}, we get 
\begin{equation}
\label{104}
\prob\left( \wedge _{\ell=0}^{k}{\cal E}_\ell \right)\geq 1-\sum_{\ell=0}^{k}\prob(\neg {\cal E}_\ell)
\geq 1/4 
\end{equation}

Since $u$ is Gaussian, $\lambda , z$ can be viewed as independent gaussian random variables. So we have that ${\cal E}_{k+1}$ is independent of ${\cal E}_\ell,\ell=0,1,2,\ldots , k$. Therefore,~\Cref{fact:3} and \eqref{104} imply that 
\begin{equation*}
\prob\left( \wedge_{\ell=0}^{k+1} {\cal E}_\ell\right)\geq  \frac{1}{40}    k^{-10/\delta^2}.
\end{equation*}
\end{proof}

For rest of the proof, we assume that the event $\wedge _{\ell=0}^{k+1}{\cal E}_\ell$ occurs. Now, for all $\ell \in [k]$, 
\begin{equation}\label{722}
u\cdot (\zeta_\ell-b) \stackrel{\eqref{eq:step2}}{\leq} z\cdot (\zeta_\ell-b)\leq
2\sqrt{\ln k} \dia(K),
\end{equation}
where the last inequality follows from the fact that the event $\cE_\ell$ has occured.
Also, since $z\perp (a-b)$:
\begin{equation}\label{723}
u\cdot (a-b)=\lambda |a-b|\geq 3\sqrt{\ln k} \dia(K),
\end{equation}
where the last inequality follows from the fact that the event $\cE_{k+1}$ has occured and $|a-b| \geq \delta \dia(K)$. 
The above two inequalities along with~\eqref{eq:step1} imply that 
$$u\cdot a - \max_{y\in K}u\cdot y \geq \lambda |a-b|/3.$$


Now, event $\cE_0$  implies that 
$$|u|\leq \lambda + |z|\leq \lambda+4\sqrt m, $$
and so, 
$$ u\cdot a - \max_{y\in K}u\cdot y  \geq \frac{\lambda \delta \dia(K) |u|}{3(\lambda + 4 \sqrt{m})},$$ 
The desired result now follows from~\Cref{cor:event} and the fact that $\lambda \cdot \delta \geq 3 \sqrt{\ln k}$ (\Cref{fact:3}). 
\end{proof}


